\subsection{Matched-support tests of information sufficiency}
The Brazilian analysis compares information sets on exactly the same footprint, endpoints, area weights and mapped domains. The native-cell masses $q_{it}$ are first summed over nested, non-overlapping square blocks with sides $s=1,2,4,\ldots,256$ native cells. All 65,536 cells remain included at every scale. This operation coarsens the available mass observations; it neither changes the cover classes nor creates observations at a different effective sensor resolution. The calculations are exploratory analyses of a fixed, geometry-selected footprint, not a probability sample or a preregistered ecological test.

For a block $b$, let $q_{bt}=\sum_{i\in b}q_{it}$. A block is group-pure only if every constituent cell is certified pure by integer overlap counts, and group-present if any constituent cell has positive-area group overlap. Conditional on observing only the block mass, the sharp endpoint set is
\begin{equation}
\widehat{\mathcal I}_{t,s}^g=
\left[\sum_{b\in P_{t,s}^g}q_{bt},\;\sum_{b\in H_{t,s}^g}q_{bt}\right].\label{eq:block}
\end{equation}
The difference set uses cross-endpoint extremes as above. Finer mass observations impose additional constraints, so the sets are nested under coarsening. Width reduction relative to the one-window observation quantifies an increment in identification information, not classification accuracy, confidence coverage or ecological predictive performance. Spatial cancellation is calculated separately from the same block-mass differences using Eq.~\ref{eq:cancel}; it describes offsets in published point maps, not independently measured biological gains and losses.

To establish attainability rather than infer meaning from wide intervals alone, two extremal allocations are constructed for each mapped group. The minimum-change witness allocates all mixed-cell mass to the group in 2015 and none in 2023; the maximum-change witness reverses those choices. Pure cells retain all their mass and absent cells none. Each witness preserves every observed native-cell total and every mapped fraction, with nonnegative complement mass. They are separate hypothetical allocations under the stated model, not estimated trajectories or jointly imposed extrema for overlapping groups. The Supplementary Information gives the construction and nesting argument.

\paragraph{Coherent nested-group allocation.}
Formations ($A$), added savanna ($S$) and their complement ($O$) form a disjoint partition; formations plus savanna is $H=A+S$. A joint allocation uses one nonnegative three-component mass partition in every cell and endpoint. Its sharp difference set in coordinates $(a,h)=(\Delta Q^A,\Delta Q^H)$ is
\begin{equation}
\mathcal D_{\rm joint}=\{(a,h):a\in\widehat{\mathcal D}^A,\ h\in\widehat{\mathcal D}^H,\ h-a\in\widehat{\mathcal D}^S\}.\label{eq:joint}
\end{equation}
The added-savanna interval is calculated from its own integer overlap counts. This joint set imposes coherence omitted by the Cartesian product of separate group intervals. The Supplementary Information establishes sharpness and describes four sign-quadrant witnesses reconstructed to preserve every native-cell mass and available component.

\subsection{Sensitivity to within-cell density contrast}
A separate sensitivity family asks how strong an additional allocation restriction must be to identify a sign. In a mixed cell with group fraction $f$, let the nonnegative component densities be $d_g$ and $d_h$. Impose $d_g\le R d_h$ and $d_h\le R d_g$ for a common assumed $R\ge1$ across mixed cells and both endpoints. With full-cell mass $q$, the sharp group-mass bounds are
\begin{equation}
q\frac{f}{f+R(1-f)}\ \le\ s_g\ \le\
q\frac{Rf}{Rf+1-f}.\label{eq:contrast}
\end{equation}
Pure and absent cells retain their topologically certified endpoint allocations. The $R=1$ limit is homogeneous allocation, whereas $R=\infty$ recovers the unrestricted primary model. No finite $R$ is inferred from the data or adopted as an ecological fact. The zero of the aggregate difference lower bound locates the loss of strict positive-sign identification within this sensitivity family. A finite bound would require independent justification before supporting a physical conclusion; this sensitivity calculation does not calibrate measurement error or supply the missing paired reference information.
